%% kapa: exam=1 round=37 year=2026 period=1 form=A subject=civil_law subject_ko=민법 seq=17 number=17
%% source: Q-Net 감정평가사 기출문제 (cst00309) artlSeq=5258924 file=raw/1008/5258924/2026년 제37회 감정평가사 1차 1교시.pdf
%% conversion: kordoc 4.24.51 -> tools/gen1.py (auto); review=proofread-sonnet; flags=none
\begin{question}{17}
민법 제145조 소정의 ‘법정추인사유’에 해당하지 않는 경우는?
\begin{choices}
  \item 취소권자의 상대방이 취소권자에게 이행을 청구한 경우
  \item 취소권자가 채무자로서 저당권을 설정한 경우
  \item 취소권자가 채권자로서 강제집행을 한 경우
  \item 취소권자가 취소할 수 있는 매매행위에 의해 취득한 소유권을 양도한 경우
  \item 취소권자인 채무자가 취소할 수 있는 행위에 의해 성립된 채무를 소멸시키고 그 대신 다른 채무를 성립시키는 계약을 체결한 경우
\end{choices}
\end{question}
\answer{1}
